% Q05_gold_answers.tex — GRADER CALIBRATION FIXTURE (correct, paraphrased).
% A drop-in replacement for paper/Q05_answers.tex (spec §4.7, §6.12): every answer is
% correct, but the order of the steps, the wording, the chosen witnesses and the
% blueprint's order are deliberately far from the reference, so the grader must
% recognise correctness by substance (the rubric's probes), not by matching text.
% B1 applies both call rules before it simplifies anything and exhibits another
% outcome; B2 exhibits a three-limb state; B3 argues the load branch first and writes
% its invariant as a difference; B4 words its fields differently. Reference-class (hook
% rule 1, DELTAS Δ125): a correct answer set, with a blueprint of Q05.H4d and H4v, so
% open it only after Q05 is done.
%
% EXPECTED[B1][recomputation][4]: correct; Probe A's Sufficient verifies with 2 <= L spliced in, and Probe A''' passes on s = [1, 1], c = [0, 6], c' = [2, 0] (results 1 and 6) for each of the two conjuncts named. No finding.
% EXPECTED[B2][recomputation][4]: correct; Probe B verifies each condition decided valid, and Probe B' passes on L = 3, i = 2, c = [0, 0, 4], x = 4, c' = [0, 8, 0]. No finding.
% EXPECTED[B3][statement-fidelity][4] [induction-soundness][4] [case-exhaustiveness][4] [blueprint-quality][4] [gap-accounting][4]: correct; Probe C filled from this blueprint (0 <= i <= |s|, Ones(s) - c == BitCount(x) + Ones(s[i..]), decreases |s| - i, x) verifies, and Probe C'' verifies with the overflow paragraph's facts. No finding.
% EXPECTED[B4][card-fidelity][4]: every field agrees with the pinned range (lines 1366–1386); Artifact correctly says the range is silent on a code release. No finding on any field.
% Expected per-problem grades: B1 = 4, B2 = 4, B3 = 4, B4 = 4 (3 is tolerated only with a declared-gap note and no flaw).
% A calibration run FAILS if any problem scores <= 2 or if any finding claims a wrong step,
% an uncovered instance or a contradiction with the range.
\documentclass[11pt]{article}
\input{preamble}

\title{Q05 --- Frames, ownership, modes, wp$\to$VC \\ \large Written answers}
\author{learner}
\date{}

\begin{document}
\maketitle

% ---------------------------------------------------------------------------
\section*{Problem B1 (hand-calculation)}\label{ans:B1}
\begin{answer}
\textbf{(a)} I name the contents $s$ (start), $c_1$ (after
\dfy{Drain(a, 0)}) and $c_2$ (after \dfy{Drain(a, 1)}), and apply the rules
first, simplifying only at the end. By the sequence rule I go backwards
through the three statements. \\
Assignment rule on \dfy{t := x + y}: $R$ becomes $x + y = s[0] + s[1] \wedge
c_2[0] = 0 \wedge c_2[1] = 0$. \\
Call rule on \dfy{y := Drain(a, 1)} (precondition at the argument $1$, then a
fresh outcome $c_2, y_1$ under the postcondition):
$F_2 \equiv 1 < L \wedge \forall c_2, y_1 .\; (y_1 = c_1[1] \wedge c_2[1] = 0
\wedge \forall j \neq 1 .\; 0 \leq j < L \impl c_2[j] = c_1[j]) \impl (x + y_1
= s[0] + s[1] \wedge c_2[0] = 0 \wedge c_2[1] = 0)$. \\
Call rule on \dfy{x := Drain(a, 0)}:
$F_1 \equiv 0 < L \wedge \forall c_1, x_1 .\; (x_1 = s[0] \wedge c_1[0] = 0
\wedge \forall j \neq 0 .\; 0 \leq j < L \impl c_1[j] = s[j]) \impl
\subst{F_2}{x}{x_1}$. \\
Simplifying, innermost first. In $F_2$ the hypotheses pin down everything the
conclusion reads: $y_1 = c_1[1]$, $c_2[1] = 0$, and $c_2[0] = c_1[0]$ (frame,
$j = 0$, which is in range because $1 < L$). Some outcome satisfies them
($c_1$ with limb 1 zeroed), so $F_2 \equiv 1 < L \wedge x + c_1[1] = s[0] +
s[1] \wedge c_1[0] = 0$. In $F_1$, the conjunct $1 < L$ does not depend on the
outcome and moves out; then the hypotheses give $x_1 = s[0]$, $c_1[0] = 0$
and, by the frame at $j = 1$, $c_1[1] = s[1]$, so the rest of the conclusion
is true. What remains is $0 < L \wedge 1 < L$, that is $2 \leq L$, whatever
$s$ holds.

\textbf{(b)} $L = 2 \impl 2 \leq L$, which holds. \\
Without the frame postcondition, the call rule's hypotheses shrink to
$y_1 = c_1[1] \wedge c_2[1] = 0$ and $x_1 = s[0] \wedge c_1[0] = 0$: nothing
ties $c_2[0]$ to $c_1[0]$ or $c_1[1]$ to $s[1]$, so $F_2$ and the wp are false
and the condition fails. The sum $t = s[0] + s[1]$ and $c_2[0] = 0$ can no
longer be proved; $c_2[1] = 0$ still can, as the second call's own
postcondition. \\
An outcome that the weaker contract allows: $s = [1, 1]$; the first call
returns $1$ and leaves $c_1 = [0, 6]$; the second returns $6$ and leaves
$c_2 = [2, 0]$. Then $t = 7 \neq 2$ and $c_2[0] = 2 \neq 0$. \\
No execution ends like this: the body of \dfy{Drain} only zeroes limb $i$, so
every run from $[1, 1]$ returns $2$ and leaves $[0, 0]$. The failure is the
contract's, not the program's: \dfy{Collect} can use only what \dfy{Drain}
promises, and the deleted line was the only promise about the other limb.

\textbf{(c)} In Verus the frame is the type: \vrs{&mut Vec<u64>} hands the
callee the whole vector, length and all, so its \vrs{forall} clause plus the
length clause is what Dafny's frame postcondition says. \\
The extra obligation is that \vrs{x + y} fits in a \vrs{u64}. The caller gets
it from its \vrs{< 1000} bounds on both start limbs, \vrs{drain}'s promise
that it returns the old limb, and the first call's \vrs{forall}, which keeps
limb 1 for the second call. \\
Deleting that \vrs{forall} makes the overflow check at \vrs{x + y} fail, and
the postcondition's sum and \vrs{final(v)@[0] == 0}; \vrs{final(v)@[1] == 0}
survives.
\end{answer}

% ---------------------------------------------------------------------------
\section*{Problem B2 (hand-calculation)}\label{ans:B2}
\begin{answer}
\textbf{(a)} With $I \equiv 0 \leq i \leq L \wedge \forall j .\; 0 \leq j < i \impl c[j] = 0$:
\begin{itemize}
\item Entry: $i \assign 0$ turns $I$ into $0 \leq L$ and a statement about no
  $j$ at all. Valid.
\item Maintenance: backwards through the body, \dfy{i := i + 1} first, by the
  assignment rule: $0 \leq i + 1 \leq L \wedge \forall j < i + 1 .\; c'[j] =
  0$; then the call, by the call rule: $i < L \wedge \forall c', x' .\;
  (x' = c[i] \wedge c'[i] = 0 \wedge \forall j \neq i .\; 0 \leq j < L \impl
  c'[j] = c[j]) \impl (\ldots)$. Under $I$ and the guard, the guard discharges
  the precondition and $i + 1 \leq L$; limbs below $i$ are kept by the frame
  and were zero; limb $i$ is zero by the postcondition. Valid.
\item Use at exit: $I$ with $i \geq L$ forces $i = L$, which is the
  postcondition. Valid.
\item Termination: while $i < L$, the measure $L - i \geq 1$; the body adds
  one to $i$ and leaves the length alone, so $L - i$ drops by one. Valid.
\end{itemize}

\textbf{(b)} Only maintenance breaks: the other three never used what
\dfy{Drain} keeps. With the outcome constrained only by $x' = c[i]$ and
$c'[i] = 0$: take $L = 3$, $i = 2$, $c = [0, 0, 4]$, and the outcome $x' = 4$,
$c' = [0, 8, 0]$. The hypotheses hold, and $c'[1] = 8$ breaks the
conclusion. \\
No run reaches it, because the real \dfy{Drain} writes one limb. And no
invariant, however strong, rescues the loop: whatever it says about the limbs
before the last call, that call may now overwrite all of them except its own.
The weak spot is the callee's postcondition, not the loop's invariant, which
is the reverse of Q03's B2, where a correct program failed because its
invariant missed a conjunct.

\textbf{(c)} Two fail: keeping \vrs{i <= v.len()} through the body, and the
\vrs{decreases}, which reads \vrs{v.len()}. Without the length clause the
caller cannot tell that the call left the length alone. A Dafny array's
length is fixed when it is created and is not something \dfy{modifies a}
lets the callee change, so \dfy{Drain} need not mention it.
\end{answer}

% ---------------------------------------------------------------------------
\section*{Problem B3 (informal-proof-first)}\label{ans:B3}
\begin{answer}
\textbf{Statement.} \dfy{CountOnes(s)} returns $\mathrm{Ones}(s)$ for every $s$;
the Verus loop, given \vrs{ones(v@) <= u64::MAX}, never overflows.

\textbf{Induction on / loop invariant / decreases measure:} invariant and
measure below; the measure is the pair $(|s| - i,\, x)$, first component
first.

\textbf{Base} (entry: the invariant holds before the first iteration)\textbf{:}
with $i = x = c = 0$, the difference $\mathrm{Ones}(s) - 0$ equals
$\mathrm{BitCount}(0) + \mathrm{Ones}(s[0..]) = 0 + \mathrm{Ones}(s)$.

\textbf{IH or invariant (stated as a formula):}
$0 \leq i \leq |s|$ and $\mathrm{Ones}(s) - c = \mathrm{BitCount}(x) +
\mathrm{Ones}(s[i..])$: what is left to count is the bits still in $x$ plus
those of the limbs not yet loaded.

\textbf{Step} (one paragraph per branch of the body: the invariant is kept,
and the measure goes down while staying bounded below)\textbf{:} \\
\emph{Loading} ($x = 0$, so $i < |s|$). Afterwards $x = s[i]$ and $i$ is one
larger. The right-hand side was $0 + \mathrm{Ones}(s[i..])$, which unfolds to
$\mathrm{BitCount}(s[i]) + \mathrm{Ones}(s[i..][1..])$, and the code's own
assertion rewrites $s[i..][1..]$ as $s[i + 1..]$: exactly the new right-hand
side; $c$ is untouched. The first component $|s| - i$ falls by one, so the
pair falls even though $x$ jumps up. \\
\emph{Counting a bit} ($x > 0$). $c$ grows by $x \bmod 2$, $x$ halves, $i$
stays. Unfolding BitCount once at a non-zero $x$ shows that the right-hand
side shrinks by exactly $x \bmod 2$, as the left does. The first component
stays and the second, $x$, falls, because halving a positive number makes it
smaller. \\
Both components are natural numbers, so the order is well founded.

\textbf{Use at exit:} the loop stops when $x = 0$ and $i \geq |s|$, so
$i = |s|$; the right-hand side is $0 + \mathrm{Ones}([]) = 0$, so $c =
\mathrm{Ones}(s)$.

\textbf{Overflow in H4v} (why \texttt{c + x \% 2} fits in a \vrs{u64})\textbf{:}
the loop carries \vrs{ones(v@) <= u64::MAX} as an invariant, since it sees
nothing else from outside. In the counting branch $x > 0$, and one unfolding
gives $\mathrm{bit\_count}(x) \geq x \bmod 2$; with the invariant,
$c + x \bmod 2 \leq c + \mathrm{bit\_count}(x) \leq \mathrm{ones}(v) \leq
2^{64} - 1$.

\textbf{Helper lemmas needed (this is the blueprint;} each by its statement:
an unfolding, the given assertion, or a new lemma\textbf{):} unfoldings only:
BitCount at a positive $x$ and at $0$, Ones at a non-empty sequence and at
$[]$; the given $s[i..][1..] = s[i + 1..]$; and $s[0..] = s$ before the loop
(in Verus an extensional equality, asserted with \vrs{=~=}, as is the
suffix fact). No new lemma.

\textbf{Gaps I know about} (declared gaps score better than silent ones): none.
\end{answer}

% ---------------------------------------------------------------------------
\section*{Problem B4 (paper card)}\label{ans:B4}
\begin{answer}
\textbf{Source:} cheatsheet entry 03.12, lines 1366--1386 (the pinned range).

\textbf{Task shape:} prove or refute one property of a small C program, a
predicate attached to a line. Lemur runs a proof calculus (eight rules over
the program, at most one assumption and a stack of goals); GPT supplies
candidate invariants and sub-goals, written as assertions at marked lines,
and the run ends in success (the property holds on all executions) or fail.

\textbf{Oracle:} an off-the-shelf C verifier (esbmc or UAutomizer) answering
True, False or Unknown, sound but not complete. A candidate is pushed only
once the verifier proves that it implies the goal, and the two soundness
theorems hold for any oracle.

\textbf{What is trusted:} the verifier's soundness and the calculus's
(Theorems 3.1 and 3.2, and the stability condition behind the case split of
Lemma 2.1, whose check the paper leaves undescribed); GPT's suggestions are
not trusted at all.

\textbf{Cheat surface:} the canonical boundary closes the obvious one, since a
bad suggestion cannot be accepted unproved. The rest: stability checking is
unspecified; the hard set was picked because the baselines fail on it and
holds only true properties; the very small programs may be contaminated.

\textbf{Metric:} counts of benchmarks verified within a time limit. Code2Inv,
133 programs at 10 minutes: GPT-3.5 103, GPT-4 107 (esbmc 68; Code2Inv's own
92, quoted at 1 hour). Hard SV-COMP, 47 programs at 15 minutes: 14--15 and
24--25 (19--20 for GPT-4 without repair), versus 1 each for UAutomizer and
esbmc, rising to 2 and 4 when given 12 hours.

\textbf{Budget:} $k$ proposals per goal (needed, with proposals at earlier
lines, for termination), a 30-second cap per verifier call, and the 10- or
15-minute limit per benchmark. No token or money figures; token limits appear
only as a weakness.

\textbf{Variance:} three runs per configuration, reported as ranges; nothing
about seeds or confidence intervals.

\textbf{Artifact:} the calculus and Algorithm 1, run with GPT-3.5 or GPT-4 over
esbmc or UAutomizer, on Code2Inv and the hard SV-COMP subset. Whether code or
prompts are released, the range does not say; one would need the paper's
availability statement.
\end{answer}

\end{document}
